\Chapter{48}

\Verse{1}
{
    \zA{话说薛蟠听见如此说了，气方渐平。 三五日后，疼痛虽愈，伤痕未平，只装病 在家，愧见亲友。
        展眼已到十月，因有各铺面伙计内有算年账要回家的，少不得家里治酒饯行。 内有一个张德辉，自幼在薛蟠当铺内揽总，家内也有了二三千金的过活，今岁也要
        回家，明春方来。 因说起：“今年纸札香料短少，明年必是贵的。 明年先打发大小 儿上来，当铺里照管，赶端阳前，我顺路就贩些纸札香扇来卖。 除去关税花销，稍
        亦可以剩得几倍利息。”}
}
{
    \eA{But to resume our story. After hearing his mother's arguments, Hsüeh P'an's
        indignation gradually abated. But notwithstanding that his pains and aches
        completely disappeared, in three or five days' time, the scars of his wounds
        were not yet healed and shamming illness, he remained at home; so ashamed
        was he to meet any of his relations or friends. In a twinkle, the tenth moon
        drew near;}
}
{
}

\Verse{2}
{
    \zA{薛蟠听了，心下忖度：“如今我捱了打正难见人，想着要 躲避一年半载又没处去躲。 天天装病，也不是常法儿。 况且我长了这么大，文不文
        武不武，虽说做买卖，究竟戥子、算盘从没拿过，地土风俗、远近道路又不知道。 不如也打点几个本钱和张德辉逛一年来，赚钱也罢不赚钱也罢，且躲躲羞去。 二则
        逛逛山水也是好的。” 心内主意已定，至酒席散后，便和气平心与张德辉说知，命 他等一二日，一同前往。}
}
{
    \eA{and as several among the partners in the various shops, with which he was
        connected, wanted to go home, after the settlement of the annual accounts,
        he had to give them a farewell spread at home. In their number was one Chang
        Te-hui, who from his early years filled the post of manager in Hsüeh P'an's
        pawnshop; and who enjoyed in his home a living of two or three thousand
        taels.}
}
{
}

\Verse{3}
{
    \zA{晚间薛蟠告诉他母亲，薛姨妈听了，虽是喜欢，但又恐他 在外生事，花了本钱倒是末事。 因此不叫他去，只说：“你好歹跟着我，我还放心 些。
        况且也不用这个买卖，等不着这几百银子使。” 薛蟠主意已定，那里肯依？ 只 说：“天天又说我不知世务，这个也不知，那个也不学； 如今我发狠把那些没要紧
        的都断了，如今要成人立事，学习买卖，又不准我了。 叫我怎么样呢？ 我又不是个 丫头，把我关在家里，何日是个了手？}
}
{
    \eA{His purpose too was to visit his native place this year, and to return the
        following spring. "Stationery and perfumery have been so scarce this year,"
        he consequently represented, "that prices will next year inevitably be high;}
}
{
}

\Verse{4}
{
    \zA{况且那张德辉又是个有年纪的，咱们和他是 世家，我同他怎么得有错？ 我就有一时半刻不好的去处，他自然说我劝我，就是东
        西贵贱行情，他是知道的，自然色色问他，何等顺利，倒不叫我去!过两日，我不 告诉家里，私自打点了走，明年发了财回来，才知道我呢！” 说毕，赌气睡觉去了。
        薛姨妈听他如此说，因和宝钗商议。 宝钗笑道：“哥哥果然要经历正事，倒也 罢了。 只是他在家里说着好听，到了外头，旧病复发，难拘束他了。 但也愁不得许
        多。}
}
{
    \eA{so when next year comes, what I'll do will be to send up my elder and
        younger sons ahead of me to look after the pawnshop, and when I start on my
        way back, before the dragon festival, I'll purchase a stock of paper, scents
        and fans and bring them for sale. And though we'll have to reduce the
        duties, payable at the barriers, and other expenses, there will still remain
        for us a considerable percentage of profit."}
}
{
}

\Verse{5}
{
    \zA{他若是真改了，是他一生的福； 若不改，妈妈也不能又有别的法子：一半尽人 力，一半听天罢了。 这么大人了，若只管怕他不知世路，出不得门，干不得事，今
        年关在家里，明年还是这个样儿。 他既说的名正言顺，妈妈就打量着丢了一千、八 百银子，竟交与他试一试。 横竖有伙计帮着他，也未必好意思哄骗他的。 二则他出
        去了，左右没了助兴的人，又没有倚仗的人，到了外头，谁还怕谁？ 有了的吃，没 了的饿着，举眼无靠，他见了这样，只怕比在家里省了事也未可知。”}
}
{
    \eA{This proposal set Hsüeh P'an musing, "With the dressing I've recently had,"
        he pondered, "I cannot very well, at present, appear before any one. Were
        the fancy to take me to get out of the way for half a year or even a year,
        there isn't a place where I can safely retire. And to sham illness, day
        after day, isn't again quite the right thing!}
}
{
}

\Verse{6}
{
    \zA{薛姨妈听了， 思忖半晌道：“倒是你说的是。 花两个钱叫他学些乖来也值。” 商议已定，一宿无 话。
        至次日，薛姨妈命人请了张德辉来在书房中，命薛蟠款待酒饭。 自己在后廊下 隔着窗子，千言万语嘱托张德辉照管照管。 张德辉满口应承，吃过饭告辞，又回说：
        “十四日是上好出行日期，大世兄即刻打点行李，雇下骡子，十四日一早就长行 了。” 薛蟠喜之不尽，将此话告诉了薛姨妈。}
}
{
    \eA{In addition to this, here I've reached this grown-up age, and yet I'm
        neither a civilian nor a soldier. It's true I call myself a merchant; but
        I've never in point of fact handled the scales or the abacus. Nor do I know
        anything about our territories, customs and manners, distances and routes.}
}
{
}

\Verse{7}
{
    \zA{薛姨妈和宝钗香菱并两个年老的嬷嬷，连日打点行装，派下薛蟠之奶公老苍头 一名，当年谙事旧仆二名，外有薛蟠随身常使小厮二名：主仆一共六人。 雇了三辆
        大车，单拉行李使物，又雇了四个长行骡子。 薛蟠自骑一匹家内养的铁青大走骡， 外备一匹坐马。 诸事完毕，薛姨妈宝钗等连夜劝戒之言，自不必备说。 至十三日，
        薛蟠先去辞了他母舅，然后过来辞了贾宅诸人，贾珍等未免又有饯行之说，也不必 细述。}
}
{
    \eA{So wouldn't it be advisable that I should also get ready some of my capital,
        and go on a tour with Chang Te-hui for a year or so? Whether I earn any
        money or not, will be equally immaterial to me. More, I shall escape from
        all disgrace. It will, secondly, be a good thing for me to see a bit of
        country."}
}
{
}

\Verse{8}
{
    \zA{至十四日一早，薛姨妈宝钗等直同薛蟠出了仪门，母女两个四只眼看他去了 方回来。 薛姨妈上京带来的家人不过四五房，并两三个老嬷嬷小丫头，今跟了薛蟠一
        去，外面只剩了一两个男子。 因此薛姨妈即日到书房，将一应陈设玩器并帘帐等物 尽行搬进来收贮，命两个跟去的男子之妻，一并也进来睡觉。 又命香菱将他屋里也
        收拾严紧，“将门锁了，晚上和我去睡。” 宝钗道：“妈妈既有这些人作伴，不如 叫菱姐姐和我作伴去。}
}
{
    \eA{This resolution once arrived at in his mind, he waited until they rose from
        the banquet, when he, with calmness and equanimity, brought his plans to
        Chang Te-hui's cognizance, and asked him to postpone his departure for a day
        or two so that they should proceed on the journey together. In the evening,
        he imparted the tidings to his mother.}
}
{
}

\Verse{9}
{
    \zA{我们园里又空，夜长了，我每夜做活，越多一个人，岂不越 好？” 薛姨妈笑道：“正是我忘了，原该叫他和你去才是。 我前日还和你哥哥说：
        文杏又小，到三不着两的； 莺儿一个人，不够伏侍的。 还要买一个丫头来你使。” 宝钗道：“买的不知底里，倘或走了眼，花了钱事小，没的淘气。
        倒是慢慢打听着， 有知道来历的，买个还罢了。” 一面说，一面命香菱收拾了衾褥妆奁，命一个老嬷 嬷并臻儿送至蘅芜院去，然后宝钗和香菱才同回园中来。}
}
{
    \eA{Mrs. Hsüeh, upon hearing his intention, was albeit delighted, tormented with
        fresh misgivings lest he should stir up trouble abroad,—for as far as the
        expense was concerned she deemed it a mere bagatelle,—and she consequently
        would not permit him to go. "You have," she reasoned with him, "to take
        proper care of me, so that I may be able to live in peace.}
}
{
}

\Verse{10}
{
    \zA{香菱向宝钗道：“我原要和太太说的，等大爷去了，我和姑娘做伴去。 我又恐 怕太太多心，说我贪着园里来玩，谁知你竟说了。” 宝钗笑道：“我知道你心里羡
        慕这园子不是一日两日的了，只是没有个空儿。 每日来一趟，慌慌张张的，也没趣 儿。 所以趁着机会，越发住上一年，我也多个做伴的，你也遂了你的心。” 香菱笑
        道：“好姑娘!趁着这个功夫，你教给我做诗罢！” 宝钗笑道：“我说你‘得陇望 蜀’呢。}
}
{
    \eA{Another thing is, that you can well dispense with all this buying and
        selling, for you are in no need of the few hundreds of taels, you may make."
        Hsüeh P'an had long ago thoroughly resolved in his mind what to do and he
        did not therefore feel disposed to listen to her remonstrances.}
}
{
}

\Verse{11}
{
    \zA{我劝你且缓一缓，今儿头一日进来，先出园东角门，从老太太起，各处各 人，你都瞧瞧，问候一声儿，也不必特意告诉他们搬进园来。 若有提起因由儿的，
        你只带口说我带了你进来做伴儿就完了。 回来进了园，再到各姑娘房里走走。” 香 菱应着，才要走时，只见平儿忙忙的走来。 香菱忙问了好，平儿只得陪笑相问。
        宝 钗因向平儿笑道：“我今儿把他带了来做伴儿，正要回你奶奶一声儿。” 平儿笑道： “姑娘说的是那里的话？ 我竟没话答言了。”}
}
{
    \eA{"You daily tax me," he pleaded, "with being ignorant of the world, with not
        knowing this, and not learning that, and now that I stir up my good
        resolution, with the idea of putting an end to all trifling, and that I wish
        to become a man, to do something for myself, and learn how to carry on
        business, you won't let me! But what would you have me do? Besides I'm not a
        girl that you should coop me up at home!}
}
{
}

\Verse{12}
{
    \zA{宝钗道：“这才是正理。 ‘店房有个 主人，庙里有个住持。’ 虽不是大事，到底告诉一声，就是园里坐更上夜的人，知 道添了他两个，也好关门候户的了。
        你回去就告诉一声罢，我不打发人说去了。” 平儿答应着，因又向香菱道：“你既来了，也不拜拜街坊去吗？” 宝钗笑道：“我 正叫他去呢。”
        平儿道：“你且不必往我们家去，二爷病了在家里呢。” 香菱答应 着去了，先从贾母处来，不在话下。}
}
{
    \eA{And when is this likely to come to an end? Chang Te-hui is, moreover, a man
        well up in years; and he is an old friend of our family, so if I go with
        him, how ever will I be able to do anything that's wrong? Should I at any
        time be guilty of any impropriety, he will be sure to speak to me, and to
        exhort me. He even knows the prices of things and customs of trade;}
}
{
}

\Verse{13}
{
    \zA{且说平儿见香菱去了，就拉宝钗悄悄说道：“姑娘可听见我们的新文没有？” 宝钗道：“我没听见新文。 因连日打发我哥哥出门，所以你们这里的事，一概不知 道；
        连姐妹们这两天没见。” 平儿笑道：“老爷把二爷打的动不得，难道姑娘就没 听见吗？” 宝钗道：“早起恍惚听见了一句，也信不真。 我也正要瞧你奶奶去呢，
        不想你来。 又是为了什么打他？” 平儿咬牙骂道：“都是那什么贾雨村，半路途中 那里来的饿不死的野杂种!认了不到十年，生了多少事出来。}
}
{
    \eA{and as I shall, as a matter of course, consult him in everything, what
        advantage won't I enjoy? But if you refuse to let me go, I'll wait for a
        couple of days, and, without breathing a word to any one at home, I'll
        furtively make my preparations and start, and, when by next year I shall
        have made my fortune and come back, you'll at length know what stuff I'm
        made off!"}
}
{
}

\Verse{14}
{
    \zA{今年春天，老爷不知 在那个地方看见几把旧扇子，回家来，看家里所有收着的这些好扇子，都不中用了， 立刻叫人各处搜求。
        谁知就有个不知死的冤家，混号儿叫做石头呆子，穷的连饭也 没的吃，偏偏他家就有二十把旧扇子，死也不肯拿出大门来。 二爷好容易烦了多少
        情，见了这个人，说之再三，他把二爷请了到他家里坐着，拿出这扇子来略瞧了一 瞧。 据二爷说，原是不能再得的，全是湘妃、棕竹、麋鹿、玉竹的，皆是古人写画
        真迹。}
}
{
    \eA{When he had done speaking, he flew into a huff and went off to sleep. Mrs.
        Hsüeh felt impelled, after the arguments she heard him propound, to
        deliberate with Pao-ch'ai. "If brother," Pao-ch'ai smilingly rejoined, "were
        in real earnest about gaining experience in some legitimate concerns, it
        would be well and good.}
}
{
}

\Verse{15}
{
    \zA{回来告诉了老爷，便叫买他的，要多少银子给他多少。 偏那石呆子说：‘我 饿死冻死，一千两银子一把，我也不卖。’ 老爷没法了，天天骂二爷没能为。 已经
        许他五百银子，先兑银子，后拿扇子，他只是不卖，只说：‘要扇子先要我的命！’ 姑娘想想，这有什么法子？ 谁知那雨村没天理的听见了，便设了法子，讹他拖欠官
        银，拿他到了衙门里去，说：‘所欠官银，变卖家产赔补。’ 把这扇子抄了来，做 了官价，送了来。 那石呆子如今不知是死是活。}
}
{
    \eA{But though he speaks, now that he is at home, in a plausible manner, the
        moment he gets abroad, his old mania will break out again, and it will be
        hard to exercise any check over him. Yet, it isn't worth the while
        distressing yourself too much about him! If he does actually mend his ways,
        it will be the happiness of our whole lives. But if he doesn't change, you
        won't, mother, be able to do anything more;}
}
{
}

\Verse{16}
{
    \zA{老爷问着二爷说：‘人家怎么弄了 来了？’ 二爷只说了一句：‘为这点子小事弄的人家倾家败产，也不算什么能为。’ 老爷听了就生了气，说二爷拿话堵老爷呢。
        这是第一件大的。 过了几日，还有几件 小的，我也记不清，所以都凑在一处，就打起来了。 也没拉倒用板子棍子，就站着，
        不知他拿什么东西打了一顿，脸上打破了两处。 我们听见姨太太这里有一种药上棒 疮的，姑娘寻一丸给我呢。” 宝钗听了，忙命莺儿去找了两丸来与平儿。}
}
{
    \eA{for though, in part, it depends on human exertion, it, in part, depends upon
        the will of heaven! If you keep on giving way to fears that, with his lack
        of worldly experience, he can't be fit to go abroad and can't be up to any
        business, and you lock him up at home this year, why next year he'll be just
        the same!}
}
{
}

\Verse{17}
{
    \zA{宝钗道： “既这样，你去替我问候罢，我就不去了。” 平儿向宝钗答应着去了，不在话下。
        且说香菱见了众人之后，吃过晚饭，宝钗等都往贾母处去了，自己便往潇湘馆 中来。 此时黛玉已好了大半了，见香菱也进园来住，自是喜欢。 香菱因笑道：“我
        这一进来了，也得空儿，好歹教给我做诗，就是我的造化了。” 黛玉笑道：“既要 学做诗，你就拜我为师。 我虽不通，大略也还教的起你。”
        香菱笑道：“果然这样， 我就拜你为师，你可不许腻烦的。”}
}
{
    \eA{Such being the case, you'd better, ma,—since his arguments are right and
        specious enough,—make up your mind to sacrifice from eight hundred to a
        thousand taels and let him have them for a try. He'll, at all events, have
        one of his partners to lend him a helping hand, one who won't either think
        it a nice thing to play any of his tricks upon him.}
}
{
}

\Verse{18}
{
    \zA{黛玉道：“什么难事，也值得去学？ 不过是起、 承、转、合，当中承、转是两副对子，平声的对仄声，虚的对实的，实的对虚的。
        若是果有了奇句，连平仄虚实不对都使得的。” 香菱笑道：“怪道我常弄本旧诗， 偷空儿看一两首，又有对的极工的，又有不对的。 又听见说，‘一三五不论，二四
        六分明。’ 看古人的诗上，亦有顺的，亦有二四六上错了的。 所以天天疑惑。 如今 听你一说，原来这些规矩，竟是没事的，只要词句新奇为上。”
        黛玉道：“正是这 个道理。}
}
{
    \eA{In the second place, there will be, when he's gone, no one to the left of
        him or to the right of him, to stand by him, and no one upon whom to rely,
        for when one goes abroad, who cares for any one else? Those who have, eat;
        and those who haven't starve.}
}
{
}

\Verse{19}
{
    \zA{词句究竟还是末事，第一是立意要紧。 若意趣真了，连词句不用修饰自是 好的，这叫做‘不以词害意’。” 香菱道：“我只爱陆放翁的‘重帘不卷留香久，
        古砚微凹聚墨多’，说的真切有趣。” 黛玉道：“断不可看这样的诗。 你们因不知 诗，所以见了这浅近的就爱，一入了这个格局，再学不出来的。 你只听我说，你若
        真心要学，我这里有《王摩诘全集》，你且把他的五言律一百首细心揣摩透熟了， 然后再读一百二十首老杜的七言律，次之再李青莲的七言绝句读一二百首。}
}
{
    \eA{When he therefore casts his eyes about him and realises that there's no one
        to depend upon, he may, upon seeing this, be up to less mischief than were
        he to stay at home; but of course, there's no saying." Mrs. Hsüeh listened
        to her, and communed within herself for a moment. "What you say is, indeed,
        right and proper!" she remarked.}
}
{
}

\Verse{20}
{
    \zA{肚子里 先有了这三个人做了底子，然后再把陶渊明、应、刘、谢、阮、庾、鲍等人的一看， 你又是这样一个极聪明伶俐的人，不用一年工夫，不愁不是诗翁了。”
        香菱听了， 笑道：“既这样，好姑娘，你就把这书给我拿出来，我带回去夜里念几首也是好的。”
        黛玉听说，便命紫鹃将王右丞的五言律拿来，递与香菱道：“你只看有红圈的，都 是我选的，有一首念一首。 不明白的问你姑娘，或者遇见我，我讲与你就是了。”}
}
{
    \eA{"And could one, by spending a small sum, make him learn something
        profitable, it will be well worth!" They then matured their plans; and
        nothing further of any note transpired during the rest of the night. The
        next day, Mrs. Hsüeh sent a messenger to invite Chang Te-hui to come round.
        On his arrival, she charged Hsüeh P'an to regale him in the library.}
}
{
}

\Verse{21}
{
    \zA{香菱拿了诗，回至蘅芜院中，诸事不管，只向灯下一首一首的读起来。 宝钗连催他 数次睡觉，他也不睡。 宝钗见他这般苦心，只得随他去了。
        一日，黛玉方梳洗完了，只见香菱笑吟吟的送了书来，又要换杜律。 黛玉笑道： “共记得多少首？” 香菱笑道：“凡红圈选的，我尽读了。” 黛玉道：“可领略了
        些没有？” 香菱笑道：“我倒领略了些，只不知是不是，说给你听听。” 黛玉笑道： “正要讲究讨论，方能长进。 你且说来我听听。”}
}
{
    \eA{Then appearing, in person, outside the window of the covered back passage,
        she made thousand of appeals to Chang Te-hui to look after her son and take
        good care of him. Chang Te-hui assented to her solicitations with profuse
        assurances, and took his leave after the collation. "The fourteenth," he
        went on to explain to Hsüeh P'an. "is a propitious day to start.}
}
{
}

\Verse{22}
{
    \zA{香菱笑道：“据我看来，诗的好 处，有口里说不出来的意思，想去却是逼真的； 又似乎无理的，想去竟是有理有情 的。”
        黛玉笑道：“这话有了些意思!但不知你从何处见得？” 香菱笑道：“我看 他《塞上》一首，内一联云：‘大漠孤烟直，长河落日圆。’ 想来烟如何直？ 日自
        然是圆的。 这‘直’字似无理，‘圆’字似太俗。 合上书一想，倒像是见了这景的。 要说再找两个字换这两个，竟再找不出两个字来。
        再还有：‘日落江湖白，潮来天 地青。’}
}
{
    \eA{So, worthy friend, you'd better be quick and pack up your baggage, and hire
        a mule, for us to begin our long journey as soon as the day dawns on the
        fourteenth." Hsüeh P'an was intensely gratified, and he communicated their
        plans to Mrs. Hsüeh. Mrs. Hsüeh then set to, and worked away, with the
        assistance of Pao-ch'ai, Hsiang Ling and two old nurses, for several
        consecutive days, before she got his luggage ready.}
}
{
}

\Verse{23}
{
    \zA{这‘白’‘青’两个字，也似无理，想来必得这两个字才形容的尽，念在 嘴里，倒像有几千斤重的一个橄榄似的。 还有‘渡头馀落日，墟里上孤烟’，这‘馀’
        字合‘上’字，难为他怎么想来!我们那年上京来，那日下晚便挽住船，岸上又没 有人，只有几棵树。 远远的几家人家作晚饭，那个烟竟是青碧连云。 谁知我昨儿晚
        上看了这两句，倒像我又到了那个地方去了。” 正说着，宝玉和探春来了，都入座听他讲诗。}
}
{
    \eA{She fixed upon the husband of Hsüeh P'an's nurse an old man with hoary head,
        two old servants with ample experience and long services, and two young
        pages, who acted as Hsüeh P'an's constant attendants, to go with him as his
        companions, so the party mustered, inclusive of master and followers, six
        persons in all. Three large carts were hired for the sole purpose of
        carrying the baggage and requisites;}
}
{
}

\Verse{24}
{
    \zA{宝玉笑道：“既是这样，也不用 看诗，‘会心处不在远’，听你说了这两句，可知三昧你已得了。” 黛玉笑道：“你
        说他这‘上孤烟’好，你还不知他这一句还是套了前人的来。 我给你这一句瞧瞧， 更比这个淡而现成。” 说着，便把陶渊明的“暧暧远人村，依依墟里烟”翻了出来，
        递给香菱。 香菱瞧了，点头叹赏，笑道：“原来‘上’字是从‘依依’两个字上化 出来的。” 宝玉大笑道：“你已得了。 不用再讲，要再讲，倒学离了。
        你就做起来 了，必是好的。”}
}
{
    \eA{and four mules, suitable for long journeys, were likewise engaged. A tall,
        dark brown, home-bred mule was selected for Hsüeh P'an's use; but a saddle
        horse, as well, was provided for him. After the various preparations had
        been effected, Mrs. Hsüeh, Pao-ch'ai and the other inmates tendered him,
        night after night, words of advice. But we can well dispense with dilating
        on this topic.}
}
{
}

\Verse{25}
{
    \zA{探春笑道：“明儿我补一个柬来，请你入社。” 香菱道：“姑娘 何苦打趣我!我不过是心里羡慕，才学这个玩罢了。” 探春黛玉都笑道：“谁不是 玩？
        难道我们是认真做诗呢!要说我们真成了诗，出了这园子，把人的牙还笑掉了 呢。” 宝玉道：“这也算自暴自弃了。 前儿我在外头和相公们商画儿，他们听见咱
        们起诗社，求我把稿子给他们瞧瞧，我就写了几首给他们看看。 谁不是真心叹服？ 他们抄了刻去了。” 探春黛玉忙问道：“这是真话么？”}
}
{
    \eA{On the arrival of the thirteenth, Hsüeh P'an went and bade good-bye to his
        maternal uncles. After which, he came and paid his farewell visit to the
        members of the Chia household. Chia Chen and the other male relatives
        unavoidably prepared an entertainment to speed him off. But to these
        festivities, there is likewise little need to allude with any minuteness.}
}
{
}

\Verse{26}
{
    \zA{宝玉笑道：“说谎的是那 架上鹦哥。” 黛玉探春听说，都道：“你真真胡闹!且别说那不成诗，便成诗，我 们的笔墨，也不该传到外头去。” 宝玉道：“这怕什么？
        古来闺阁中笔墨不要传出 去，如今也没人知道呢。” 说着，只见惜春打发了入画来请宝玉，宝玉方去了。
        香菱又逼着换出杜律，又央黛玉探春二人：“出个题目让我诌去，诌了来替我 改正。” 黛玉道：“昨夜的月最好，我正要诌一首未诌成。 你就做一首来。 ‘十四
        寒’的韵，由你爱用那几个字去。”}
}
{
    \eA{On the fourteenth, at break of day, Mrs. Hsüeh, Pao-ch'ai and the other
        members of the family accompanied Hsüeh P'an beyond the ceremonial gate.
        Here his mother and her daughter stood and watched him, their four eyes
        fixed intently on him, until he got out of sight, when they, at length,
        retraced their footsteps into the house.}
}
{
}

\Verse{27}
{
    \zA{香菱听了，喜的拿着诗回来，又苦思一回，做 两句诗； 又舍不得杜诗，又读两首：如此茶饭无心，坐卧不定。 宝钗道：“何苦自 寻烦恼？
        都是颦儿引的你，我和他算帐去!你本来呆头呆脑的，再添上这个，越发弄 成个呆子了。” 香菱笑道：“好姑娘，别混我。” 一面说，一面做了一首。 先给宝
        钗看了，笑道：“这个不好，不是这个做法。 你别害臊，只管拿了给他瞧去，看是 他怎么说。” 香菱听了，便拿了诗找黛玉。}
}
{
    \eA{Mrs. Hsüeh had, in coming up to the capital, only brought four or five
        family domestics and two or three old matrons and waiting-maids with her,
        so, after the departure on the recent occasion, of those, who followed Hsüeh
        P'an, no more than one or two men-servants remained in the outer quarters.}
}
{
}

\Verse{28}
{
    \zA{黛玉看时，只见写道是： 月桂中天夜色寒，清光皎皎影团团。 诗人助兴常思玩，野客添愁不忍观。 翡翠楼边悬玉镜，珍珠帘外挂冰盘。
        良宵何用烧银烛，晴彩辉煌映画栏。 黛玉笑道：“意思却有，只是措词不雅。 皆因你看的诗少，被他缚住了。 把这首诗 丢开，再做一首。 只管放开胆子去做。”
        香菱听了，默默的回来，越发连房也不进去，只在池边树下。 或坐在山石上出 神，或蹲在地下抠地，来往的人都诧异。}
}
{
    \eA{Mrs. Hsüeh repaired therefore on the very same day into the study, and had
        the various ornaments, bric-à-brac, curtains and other articles removed into
        the inner compound and put away.}
}
{
}

\Verse{29}
{
    \zA{李纨、宝钗、探春、宝玉等听得此言，都 远远的站在山坡上瞧着他笑。 只见他皱一回眉，又自己含笑一回。 宝钗笑道：“这 个人定是疯了。
        昨夜嘟嘟哝哝，直闹到五更才睡下。 没一顿饭的工夫，天就亮了， 我就听见他起来了，忙忙碌碌梳了头，就找颦儿去。 一回来了，呆了一天，做了一
        首又不好，自然这会子另做呢。” 宝玉笑道：“这正是‘地灵人杰’，老天生人， 再不虚赋情性的。 我们成日叹说：可惜他这么个人，竟俗了。 谁知到底有今日!可
        见天地至公。”}
}
{
    \eA{Then bidding the wives of the two male attendants, who had gone with Hsüeh
        P'an, likewise move their quarters inside, along with the other women, she
        went on to impress upon Hsiang Ling to put everything carefully away in her
        own room as well, and to lock the doors; "for," (she said), "you must come
        at night and sleep with me."}
}
{
}

\Verse{30}
{
    \zA{宝钗听了，笑道：“你能够像他这苦心就好了，学什么有个不成的 吗？” 宝玉不答。 只见香菱兴兴头头的，又往黛玉那边来了。 探春笑道：“咱们跟了去，看他有
        些意思没有。” 说着，一齐都往潇湘馆来。 只见黛玉正拿着诗和他讲究呢。 众人因 问黛玉：“做的如何？” 黛玉道：“自然算难为他了，只是还不好。
        这一首过于穿 凿了，还得另做。” 众人因要诗看时，只见做道是： 非银非水映窗寒，试看晴空护玉盘。 淡淡梅花香欲染，丝丝柳带露初干。}
}
{
    \eA{"Since you've got all these people to keep you company, ma," Pao-ch'ai
        remarked, "wouldn't it be as well to tell sister Ling to come and be my
        companion? Our garden is besides quite empty and the nights are so long! And
        as I work away every night, won't it be better for me to have an extra
        person with me?" "Quite so!" smiled Mrs. Hsüeh, "I forgot that! I should
        have told her to go with you; it's but right.}
}
{
}

\Verse{31}
{
    \zA{只疑残粉涂金砌，恍若轻霜抹玉栏。 梦醒西楼人迹绝，馀容犹可隔帘看。 宝钗笑道：“不像吟月了，月字底下添一个‘色’字，倒还使得。 你看句句倒 像是月色。
        ――也罢了，原是‘诗从胡说来’，再迟几天就好了。” 香菱自为这首 诗妙绝，听如此说，自己又扫了兴，不肯丢开手，便要思索起来。 因见他姐妹们说
        笑，便自己走至阶下竹前，挖心搜胆的，耳不旁听，目不别视。 一时探春隔窗笑说 道：“菱姑娘，你闲闲罢。”}
}
{
    \eA{It was only the other day that I mentioned to your brother that: 'Wen Hsing
        too was young, and not fit to attend to everything that turns up, that Ying
        Erh could not alone do all the waiting, and that it was necessary to
        purchase another girl for your service.'" "If we buy one, we won't know what
        she's really like!" Pao-ch'ai demurred.}
}
{
}

\Verse{32}
{
    \zA{香菱怔怔答道：“‘闲’字是‘十五删’的，错了韵 了。” 众人听了，不觉大笑起来。 宝钗道：“可真诗魔了!都是颦儿引的他！” 黛
        玉笑道：“圣人说：‘诲人不倦。’ 他又来问我，我岂有不说的理！” 李纨笑道： “咱们拉了他往四姑娘屋里去，引他瞧瞧画儿，叫他醒一醒才好。” 说着，真个出
        来拉他过藕香榭，至暖香坞中。 惜春正乏倦，在床上歪着睡午觉，画缯立在壁间， 用纱罩着。 众人唤醒了惜春，揭纱看时，十停方有了三停。}
}
{
    \eA{"If she gives us the slip, the money we may have spent on her will be a mere
        trifle, so long as she hasn't been up to any pranks! So let's quietly make
        inquiries, and, when we find one with well-known antecedents, we can
        purchase her, and, we'll be on the safe side then!" While speaking, she told
        Hsiang Ling to collect her bedding and clothes;}
}
{
}

\Verse{33}
{
    \zA{见画上有几个美人，因 指香菱道：“凡会做诗的，都画在上头，你快学罢。” 说着，玩笑了一回，各自散 去。
        香菱满心中正是想诗，至晚间，对灯出了一回神，至三更以后，上床躺下，两 眼睁睁直到五更，方才蒙？ 睡着了。 一时天亮，宝钗醒了。 听了一听，他安稳睡了，
        心下想：“他翻腾了一夜，不知可做成了？ 这会子乏了，且别叫他。” 正想着，只 见香菱从梦中笑道：“可是有了!难道这一首还不好吗？”}
}
{
    \eA{and desiring an old matron and Ch'in Erh to take them over to the Heng Wu
        Yüan, Pao-ch'ai returned at last into the garden in company with Hsiang
        Ling. "I meant to have proposed to my lady," Hsiang Ling said to Pao-ch'ai,
        "that, when master left, I should be your companion, miss;}
}
{
}

\Verse{34}
{
    \zA{宝钗听了又是可叹又是 可笑，连忙叫醒了他，问他：“得了什么？ 你这诚心都通了仙了。 学不成诗，弄出 病来呢！” 一面说，一面梳洗了，和姐妹往贾母处来。
        原来香菱苦志学诗，精血诚聚，日间不能做出，忽于梦中得了八句。 梳洗已毕， 便忙写出，来到沁芳亭。 只见李纨与众姐妹方从王夫人处回来，宝钗正告诉他们，
        说他梦中做诗说梦话，众人正笑。 抬头见他来了，就都争着要诗看。 要知端底，且看下回分解。 (本章完)}
}
{
    \eA{but I feared lest her ladyship should, with that suspicious mind of hers,
        have maintained that I was longing to come into the garden to romp. But
        who'd have thought it, it was you, after all, who spoke to her about it!"}
}
{
}

\Verse{35}
{
    \zA{}
}
{
    \eA{"I am well aware," Pao-ch'ai smiled, "that you've been inwardly yearning for
        this garden, and that not for a day or two, but with the little time you can
        call your own, you would find it no fun, were you even able to run over once
        in a day, so long as you have to do it in a hurry-scurry! Seize therefore
        this opportunity of staying, better still, for a year; as I, on my side,
        will then have an extra companion; and you, on yours, will be able to
        accomplish your wishes." "My dear miss!" laughingly observed Hsiang Ling,
        "do let's make the best of this time, and teach me how to write verses!" "I
        say," Pao-ch'ai laughed, "'you no sooner, get the Lung state than you long
        for the Shu'! I advise you to wait a bit.}
}
{
}

\Verse{36}
{
    \zA{}
}
{
    \eA{This is the first day that you spend in here, and you should, first and
        foremost, go out of the garden by the eastern side gate and look up and
        salute every one in her respective quarters commencing from our old lady.
        But you needn't make it a point of telling them that you've moved into the
        garden. If anyone does allude to the reason why you've shifted your
        quarters, you can simply explain cursorily that I've brought you in as a
        companion, and then drop the subject. On your return by and bye into the
        garden, you can pay a visit to the apartments of each of the young ladies."
        Hsiang Ling signified her acquiescence, and was about to start when she saw
        P'ing Erh rush in with hurried step.}
}
{
}

\Verse{37}
{
    \zA{}
}
{
    \eA{Hsiang Ling hastened to ask after her health, and P'ing Erh felt compelled
        to return her smile, and reciprocate her inquiry. "I've brought her in
        to-day," Pao-ch'ai thereupon smilingly said to P'ing Erh, "to make a
        companion of her. She was just on the point of going to tell your lady about
        it!" "What is this that you're saying, Miss?" P'ing Erh rejoined, with a
        smile. "I really am at a loss what reply to make to you!" "It's the right
        thing!" Pao-ch'ai answered. "' In a house, there's the master, and in a
        temple there's the chief priest.' It's true, it's no important concern, but
        something must, in fact, be mentioned, so that those, who sit up on night
        duty in the garden, may be aware that these two have been added to my rooms,
        and know when to close the gates and when to wait.}
}
{
}

\Verse{38}
{
    \zA{}
}
{
    \eA{When you get back therefore do mention it, so that I mayn't have to send
        some one to tell them." P'ing Erh promised to carry out her wishes. "As
        you're moved in here," she said to Hsiang Ling, "won't you go and pay your
        respects to your neighbours?" "I had just this very moment," Pao-ch'ai
        smiled, "told her to go and do so." "You needn't however go to our house,"
        P'ing Erh remarked, "our Mr. Secundus is laid up at home." Hsiang Ling
        assented and went off, passing first and foremost by dowager lady Chia's
        apartments. But without devoting any of our attention to her, we will revert
        to P'ing Erh. Seeing Hsiang Ling walk out of the room, she drew Pao-ch'ai
        near her. "Miss! have you heard our news?" she inquired in a low tone of
        voice.}
}
{
}

\Verse{39}
{
    \zA{}
}
{
    \eA{"I haven't heard any news," Pao-ch'ai responded. "We've been daily so busy
        in getting my brother's things ready for his voyage abroad, that we know
        nothing whatever of any of your affairs in here. I haven't even seen
        anything of my female cousins these last two days." "Our master, Mr. Chia
        She, has beaten our Mr. Secundus to such a degree that he can't budge,"
        P'ing Erh smiled. "But is it likely, miss, that you've heard nothing about
        it?" "This morning," Pao-ch'ai said by way of reply, "I heard a vague report
        on the subject, but I didn't believe it could be true. I was just about to
        go and look up your mistress, when you unexpectedly arrived. But why did he
        beat him again?" P'ing Erh set her teeth to and gave way to abuse. "It's all
        on account of some Chia Yü-ts'un or other;}
}
{
}

\Verse{40}
{
    \zA{}
}
{
    \eA{a starved and half-dead boorish bastard, who went yonder quite unexpectedly.
        It isn't yet ten years, since we've known him, and he has been the cause of
        ever so much trouble! In the spring of this year, Mr. Chia She saw somewhere
        or other, I can't tell where, a lot of antique fans; so, when on his return
        home, he noticed that the fine fans stored away in the house, were all of no
        use, he at once directed servants to go everywhere and hunt up some like
        those he had seen.}
}
{
}

\Verse{41}
{
    \zA{}
}
{
    \eA{Who'd have anticipated it, they came across a reckless creature of
        retribution, dubbed by common consent the 'stone fool,' who though so poor
        as to not even have any rice to put to his mouth, happened to have at home
        twenty antique fans. But these he utterly refused to take out of his main
        door. Our Mr. Secundus had thus a precious lot of bother to ask ever so many
        favours of people. But when he got to see the man, he made endless appeals
        to him before he could get him to invite him to go and sit in his house;
        when producing the fans, he allowed him to have a short inspection of them.
        From what our Mr. Secundus says, it would be really difficult to get any the
        like of them.}
}
{
}

\Verse{42}
{
    \zA{}
}
{
    \eA{They're made entirely of spotted black bamboo, and the stags and jadelike
        clusters of bamboo on them are the genuine pictures, drawn by men of olden
        times. When he got back, he explained these things to Mr. Chia She, who
        readily asked him to buy them, and give the man his own price for them. The
        'stone fool,' however, refused. 'Were I even to be dying from hunger,' he
        said, 'or perishing from frostbites, and so much as a thousand taels were
        offered me for each single fan, I wouldn't part with them.' Mr. Chia She
        could do nothing, but day after day he abused our Mr. Secundus as a
        good-for-nothing. Yet he had long ago promised the man five hundred taels,
        payable cash down in advance, before delivery of the fans, but he would not
        sell them.}
}
{
}

\Verse{43}
{
    \zA{}
}
{
    \eA{'If you want the fans,' he had answered, 'you must first of all take my
        life.' Now, miss, do consider what was to be done? But, Yü-ts'un is, as it
        happens, a man with no regard for divine justice. Well, when he came to hear
        of it, he at once devised a plan to lay hold of these fans, so fabricating
        the charge against him of letting a government debt drag on without payment,
        he had him arrested and brought before him in the Yamên; when he adjudicated
        that his family property should be converted into money to make up the
        amount due to the public chest;}
}
{
}

\Verse{44}
{
    \zA{}
}
{
    \eA{and, confiscating the fans in question, he set an official value on them and
        sent them over here. And as for that 'stone fool,' no one now has the
        faintest idea whether he be dead or alive. Mr. Chia She, however, taunted
        Mr. Secundus. 'How is it,' he said, 'that other people can manage to get
        them?' Our master simply rejoined 'that to bring ruin upon a person in such
        a trivial matter could not be accounted ability.' But, at these words, his
        father suddenly rushed into a fury, and averred that Mr. Secundus had said
        things to gag his mouth. This was the main cause. But several minor matters,
        which I can't even recollect, also occurred during these last few days. So,
        when all these things accumulated, he set to work and gave him a sound
        thrashing.}
}
{
}

\Verse{45}
{
    \zA{}
}
{
    \eA{He didn't, however, drag him down and strike him with a rattan or cane, but
        recklessly assaulted him, while he stood before him, with something or
        other, which he laid hold of, and broke his face open in two places. We
        understand that Mrs. Hsüeh has in here some medicine or other for applying
        on wounds, so do try, miss, and find a ball of it and let me have it!"
        Hearing this, Pao-ch'ai speedily directed Ying Erh to go and look for some,
        and, on discovering two balls of it, she brought them over and handed them
        to P'ing Erh. "Such being the case," Pao-ch'ai said, "do make, on your
        return, the usual inquiries for me, and I won't then need to go."}
}
{
}

\Verse{46}
{
    \zA{}
}
{
    \eA{P'ing Erh turned towards Pao-ch'ai, and expressed her readiness to execute
        her commission, after which she betook herself home, where we will leave her
        without further notice. After Hsiang Ling, for we will take up the thread of
        our narrative with her, completed her visits to the various inmates, she had
        her evening meal. Then when Pao-ch'ai and every one else went to dowager
        lady Chia's quarters, she came into the Hsiao Hsiang lodge. By this time
        Tai-yü had got considerably better. Upon hearing that Hsiang Ling had also
        moved into the garden, she, needless to say, was filled with delight. "Now,
        that I've come in here," Hsiang Ling then smiled and said, "do please teach
        me, at your leisure, how to write verses.}
}
{
}

\Verse{47}
{
    \zA{}
}
{
    \eA{It will be a bit of good luck for me if you do." "Since you're anxious to
        learn how to versify," Tai-yü answered with a smile, "you'd better
        acknowledge me as your tutor; for though I'm not a good hand at poetry, yet
        I know, after all, enough to be able to teach you." "Of course you do!"
        Hsiang Ling laughingly remarked. "I'll readily treat you as my tutor. But
        you mustn't put yourself to any trouble!" "Is there anything so difficult
        about this," Tai-yü pursued, "as to make it necessary to go in for any
        study? Why, it's purely and simply a matter of openings, elucidations,
        embellishments and conclusions. The elucidations and embellishments, which
        come in the centre, should form two antithetical sentences, the even tones
        must pair with the uneven. Empty words must correspond with full words;}
}
{
}

\Verse{48}
{
    \zA{}
}
{
    \eA{and full words with empty words. In the event of any out-of-the-way lines,
        it won't matter if the even and uneven tones, and the empty and full words
        do not pair." "Strange though it may appear," smiled Hsiang Ling, "I often
        handle books with old poems, and read one or two stanzas, whenever I can
        steal the time; and some among these I find pair most skilfully, while
        others don't. I have also heard that the first, third and fifth lines are of
        no consequence; and that the second, fourth and sixth must be clearly
        distinguished. But I notice that there are in the poetical works of ancient
        writers both those which accord with the rules, as well as those whose
        second, fourth and sixth lines are not in compliance with any rule. Hence it
        is that my mind has daily been full of doubts.}
}
{
}

\Verse{49}
{
    \zA{}
}
{
    \eA{But after the hints you've given me, I really see that all these formulas
        are of no account, and that the main requirement is originality of diction."
        "Yes, that's just the principle that holds good," Tai-yü answered. "But
        diction is, after all, a last consideration. The first and foremost thing is
        the choice of proper sentiments; for when the sentiments are correct,
        there'll even be no need to polish the diction; it's certain to be elegant.
        This is called versifying without letting the diction affect the
        sentiments." "What I admire," Hsiang Ling proceeded with a smile; "are the
        lines by old Lu Fang; "The double portière, when not raised, retains the
        fragrance long. An old inkslab, with a slight hole, collects plenty of ink.
        "Their language is so clear that it's charming."}
}
{
}

\Verse{50}
{
    \zA{}
}
{
    \eA{"You must on no account," Tai-yü observed, "read poetry of the kind. It's
        because you people don't know what verses mean that you, no sooner read any
        shallow lines like these, than they take your fancy. But when once you get
        into this sort of style, it's impossible to get out of it. Mark my words! If
        you are in earnest about learning, I've got here Wang Mo-chieh's complete
        collection; so you'd better take his one hundred stanzas, written in the
        pentameter rule of versification, and carefully study them, until you
        apprehend them thoroughly. Afterwards, look over the one hundred and twenty
        stanzas of Lao T'u, in the heptameter rule; and next read a hundred or two
        hundred of the heptameter four-lined stanzas by Li Ch'ing-lieu.}
}
{
}

\Verse{51}
{
    \zA{}
}
{
    \eA{When you have, as a first step, digested these three authors, and made them
        your foundation, you can take T'ao Yuan-ming, Ying, Liu, Hsieh, Yüan, Yü,
        Pao and other writers and go through them once. And with those sharp and
        quick wits of yours, I've no doubt but that you will become a regular poet
        before a year's time." "Well, in that case," Hsiang Ling smiled, after
        listening to her, "bring me the book, my dear miss, so that I may take it
        along. It will be a good thing if I can manage to read several stanzas at
        night." At these words, Tai-yü bade Tzu Chüan fetch Wang Tso-ch'eng's
        pentameter stanzas. When brought, she handed them to Hsiang Ling. "Only
        peruse those marked with red circles" she said. "They've all been selected
        by me.}
}
{
}

\Verse{52}
{
    \zA{}
}
{
    \eA{Read each one of them; and should there be any you can't fathom, ask your
        miss about them. Or when you come across me, I can explain them to you."
        Hsiang Ling took the poems and repaired back to the Heng Wu-yüan. And
        without worrying her mind about anything she approached the lamp and began
        to con stanza after stanza. Pao-ch'ai pressed her, several consecutive
        times, to go to bed; but as even rest was far from her thoughts, Pao-ch'ai
        let her, when she perceived what trouble she was taking over her task, have
        her own way in the matter. Tai-yü had one day just finished combing her hair
        and performing her ablutions, when she espied Hsiang Ling come with smiles
        playing about her lips, to return her the book and to ask her to let her
        have T'u's poetical compositions in exchange.}
}
{
}

\Verse{53}
{
    \zA{}
}
{
    \eA{"Of all these, how many stanzas can you recollect?" Tai-yü asked, smiling.
        "I've read every one of those marked with a red circle," Hsiang Ling
        laughingly rejoined. "Have you caught the ideas of any of them, yes or no?"
        Tai-yü inquired. "Yes, I've caught some!" Hsiang Ling smiled. "But whether
        rightly or not I don't know. Let me tell you." "You must really," Tai-yü
        laughingly remarked, "minutely solicit people's opinions if you want to make
        any progress. But go on and let me hear you." "From all I can see," Hsiang
        Ling smiled, "the beauty of poetry lies in certain ideas, which though not
        quite expressible in words are, nevertheless, found, on reflection, to be
        absolutely correct.}
}
{
}

\Verse{54}
{
    \zA{}
}
{
    \eA{Some may have the semblance of being totally devoid of sense, but, on second
        thought, they'll truly be seen to be full of sense and feeling." "There's a
        good deal of right in what you say," Tai-yü observed. "But I wonder how you
        arrived at this conclusion?" "I notice in that stanza on 'the borderland,'
        the antithetical couplet:  "In the vast desert reigns but upright mist. In
        the long river setteth the round sun. "Consider now how ever can mist be
        upright? The sun is, of course, round. But the word 'upright' would seem to
        be devoid of common sense; and 'round' appears far too commonplace a word.
        But upon throwing the whole passage together, and pondering over it, one
        fancies having seen the scenery alluded to.}
}
{
}

\Verse{55}
{
    \zA{}
}
{
    \eA{Now were any one to suggest that two other characters should be substituted
        for these two, one would verily be hard pressed to find any other two as
        suitable. Besides this, there's also the couplet:  "When the sun sets,
        rivers and lakes are white; When the mist falls, the heavens and earth
        azure. "Both 'white' and 'azure', apparently too lack any sense; but
        reflection will show that these two words are absolutely necessary to bring
        out thoroughly the aspect of the scenery. And in conning them over, one
        feels just as if one had an olive, weighing several thousands of catties, in
        one's mouth, so much relish does one derive from them. But there's this too:
        "At the ferry stays the setting sun,  O'er the mart hangs the lonesome mist.}
}
{
}

\Verse{56}
{
    \zA{}
}
{
    \eA{"And how much trouble must these words 'stay,' and 'over, have caused the
        author in their conception! When the boats made fast, in the evening of a
        certain day of that year in which we came up to the capital, the banks were
        without a trace of human beings; and there were only just a few trees about;
        in the distance loomed the houses of several families engaged in preparing
        their evening meal, and the mist was, in fact, azure like jade, and
        connected like clouds. So, when I, as it happened, read this couplet last
        night, it actually seemed to me as if I had come again to that spot!" But in
        the course of their colloquy, Pao-yü and T'an Ch'un arrived;}
}
{
}

\Verse{57}
{
    \zA{}
}
{
    \eA{and entering the room, they seated themselves, and lent an ear to her
        arguments on the verses. "Seeing that you know so much," Pao-yü remarked
        with a smiling face, "you can dispense with reading poetical works, for
        you're not far off from proficiency. To hear you expatiate on these two
        lines, makes it evident to my mind that you've even got at their secret
        meaning." "You say," argued Tai-yü with a significant smile, "that the line:
        "'O'er (the mart) hangs the lonesome mist,' "is good; but aren't you yet
        aware that this is only plagiarised from an ancient writer? But I'll show
        you the line I'm telling you of. You'll find it far plainer and clearer than
        this." While uttering these words, she turned up T'ao Yüan-ming's,  Dim in
        the distance lies a country place;}
}
{
}

\Verse{58}
{
    \zA{}
}
{
    \eA{Faint in the hamlet-market hangs the mist; and handed it to Hsiang Ling.
        Hsiang Ling perused it, and, nodding her head, she eulogised it. "Really,"
        she smiled, the word 'over' is educed from the two characters implying
        'faint.' Pao-yü burst out into a loud fit of exultant laughter. "You've
        already got it!" he cried. "There's no need of explaining anything more to
        you! Any further explanations will, in lieu of benefiting you, make you
        unlearn what you've learnt. Were you therefore to, at once, set to work, and
        versify, your lines are bound to be good." "To-morrow," observed T'an ch'un
        with a smile; "I'll stand an extra treat and invite you to join the
        society." "Why make a fool of me, miss?" Hsiang Ling laughingly ejaculated.}
}
{
}

\Verse{59}
{
    \zA{}
}
{
    \eA{"It's merely that mania of mine that made me apply my mind to this subject
        at all; just for fun and no other reason." T'an Ch'un and Tai-yü both
        smiled. "Who doesn't go in for these things for fun?" they asked. "Is it
        likely that we improvise verses in real earnest? Why, if any one treated our
        verses as genuine verses, and took them outside this garden, people would
        have such a hearty laugh at our expense that their very teeth would drop."
        "This is again self-violence and self-abasement!" Pao-yü interposed.}
}
{
}

\Verse{60}
{
    \zA{}
}
{
    \eA{The other day, I was outside the garden, consulting with the gentlemen about
        paintings, and, when they came to hear that we had started a poetical
        society, they begged of me to let them have the rough copies to read. So I
        wrote out several stanzas, and gave them to them to look over, and who did
        not praise them with all sincerity? They even copied them and took them to
        have the blocks cut." "Are you speaking the truth?" T'an Ch'un and Tai-yü
        eagerly inquired. "If I'm telling a lie," Pao-yü laughed, "I'm like that
        cockatoo on that frame!" "You verily do foolish things!" Tai-yü and T'an
        Ch'un exclaimed with one voice, at these words. "But not to mention that
        they were doggerel lines, had they even been anything like what verses
        should be, our writings shouldn't have been hawked about outside."}
}
{
}

\Verse{61}
{
    \zA{}
}
{
    \eA{"What's there to fear?" Pao-yü smiled. "Hadn't the writings of women of old
        been handed outside the limits of the inner chambers, why, there would, at
        present, be no one with any idea of their very existence." While he passed
        this remark, they saw Ju Hua arrive from Hsi Ch'un's quarters to ask Pao-yü
        to go over; and Pao-yü eventually took his departure. Hsiang Ling then
        pressed (Tai-yü) to give her T'u's poems. "Do choose some theme," she also
        asked Tai-yü and T'an Ch'un, "and let me go and write on it. When I've done,
        I'll bring it for you to correct." "Last night," Tai-yü observed, "the moon
        was so magnificent, that I meant to improvise a stanza on it; but as I
        haven't done yet, go at once and write one using the fourteenth rhyme,
        'han,' (cool). You're at liberty to make use of whatever words you fancy."}
}
{
}

\Verse{62}
{
    \zA{}
}
{
    \eA{Hearing this, Hsiang Ling was simply delighted, and taking the poems, she
        went back. After considerable exertion, she succeeded in devising a couplet,
        but so little able was she to tear herself away from the 'T'u' poems, that
        she perused another couple of stanzas, until she had no inclination for
        either tea or food, and she felt in an unsettled mood, try though she did to
        sit or recline. "Why," Pao-ch'ai remonstrated, "do you bring such trouble
        upon yourself? It's that P'in Erh, who has led you on to it! But I'll settle
        accounts with her! You've all along been a thick-headed fool; but now that
        you've burdened yourself with all this, you've become a greater fool."
        "Miss," smiled Hsiang Ling, "don't confuse me."}
}
{
}

\Verse{63}
{
    \zA{}
}
{
    \eA{So saying, she set to work and put together a stanza, which she first and
        foremost handed to Pao-ch'ai to look over. "This isn't good!" Pao-ch'ai
        smilingly said. "This isn't the way to do it! Don't fear of losing face, but
        take it and give it to her to peruse. We'll see what she says." At this
        suggestion, Hsiang Ling forthwith went with her verses in search of Tai-yü.
        When Tai-yü came to read them, she found their text to be:  The night grows
        cool, what time Selene reacheth the mid-heavens. Her radiance pure shineth
        around with such a spotless sheen. Bards oft for inspiration raise on her
        their thoughts and eyes. The rustic daren't see her, so fears he to enhance
        his grief. Jade mirrors are suspended near the tower of malachite.}
}
{
}

\Verse{64}
{
    \zA{}
}
{
    \eA{An icelike plate dangles outside the gem-laden portière. The eve is fine, so
        why need any silvery candles burn? A clear light shines with dazzling lustre
        on the painted rails. "There's a good deal of spirit in them," Tai-yü
        smiled, "but the language is not elegant. It's because you've only read a
        few poetical works that you labour under restraint. Now put this stanza
        aside and write another. Pluck up your courage and go and work away." After
        listening to her advice, Hsiang Ling quietly wended her way back, but so
        much the more (preoccupied) was she in her mind that she did not even enter
        the house, but remaining under the trees, planted by the side of the pond,
        she either seated herself on a rock and plunged in a reverie, or squatted
        down and dug the ground, to the astonishment of all those, who went
        backwards and forwards.}
}
{
}

\Verse{65}
{
    \zA{}
}
{
    \eA{Li wan, Pao-ch'ai, T'an Ch'un, Pao-yü and some others heard about her; and,
        taking their position some way off on the mound, they watched her, much
        amused. At one time, they saw her pucker up her eyebrows; and at another
        smile to herself. "That girl must certainly be cracked!" Pao-ch'ai laughed.
        "Last night she kept on muttering away straight up to the fifth watch, when
        she at last turned in. But shortly, daylight broke, and I heard her get up
        and comb her hair, all in a hurry, and rush after P'in Erh. In a while,
        however, she returned; and, after acting like an idiot the whole day, she
        managed to put together a stanza. But it wasn't after all, good, so she's,
        of course, now trying to devise another."}
}
{
}

\Verse{66}
{
    \zA{}
}
{
    \eA{"This indeed shows," Pao-yü laughingly remarked, "that the earth is
        spiritual, that man is intelligent, and that heaven does not in the creation
        of human beings bestow on them natural gifts to no purpose. We've been
        sighing and lamenting that it was a pity that such a one as she, should,
        really, be so unpolished; but who could ever have anticipated that things
        would, in the long run, reach the present pass? This is a clear sign that
        heaven and earth are most equitable!" "If only," smiled Pao-ch'ai, at these
        words, "you could be as painstaking as she is, what a good thing it would
        be. And would you fail to attain success in anything you might take up?"
        Pao-yü made no reply.}
}
{
}

\Verse{67}
{
    \zA{}
}
{
    \eA{But realising that Hsiang Ling had crossed over in high spirits to find
        Tai-yü again, T'an Ch'un laughed and suggested, "Let's follow her there, and
        see whether her composition is any good." At this proposal, they came in a
        body to the Hsiao Hsiang lodge. Here they discovered Tai-yü holding the
        verses and explaining various things to her. "What are they like?" they all
        thereupon inquired of Tai-yü. "This is naturally a hard job for her!" Tai-yü
        rejoined. "They're not yet as good as they should be. This stanza is far too
        forced; you must write another." One and all however expressed a desire to
        look over the verses. On perusal, they read:  'Tis not silver, neither water
        that on the windows shines so cold. Selene, mark! covers, like a jade
        platter, the clear vault of heaven.}
}
{
}

\Verse{68}
{
    \zA{}
}
{
    \eA{What time the fragrance faint of the plum bloom is fain to tinge the  air,
        The dew-bedecked silken willow trees begin to lose their leaves. 'Tis the
        remains of powder which methinks besmear the golden steps. Her lustrous rays
        enshroud like light hoar-frost the jadelike  balustrade. When from my dreams
        I wake, in the west tower, all human trace is  gone. Her slanting orb can
        yet clearly be seen across the bamboo screen. "It doesn't sound like a song
        on the moon," Pao-ch'ai smilingly observed. "Yet were, after the word
        'moon', that of 'light' supplied, it would be better; for, just see, if each
        of these lines treated of the moonlight, they would be all right. But poetry
        primarily springs from nonsensical language. In a few days longer, you'll be
        able to do well."}
}
{
}

\Verse{69}
{
    \zA{}
}
{
    \eA{Hsiang Ling had flattered herself that this last stanza was perfect, and the
        criticisms, that fell on her ear, damped her spirits again. She was not
        however disposed to relax in her endeavours, but felt eager to commune with
        her own thoughts, so when she perceived the young ladies chatting and
        laughing, she betook herself all alone to the bamboo-grove at the foot of
        the steps; where she racked her brain, and ransacked her mind with such
        intentness that her ears were deaf to everything around her and her eyes
        blind to everything beyond her task. "Miss Ling," T'an Ch'un presently
        cried, smiling from inside the window, "do have a rest!" "The character
        'rest;'" Hsiang Ling nervously replied, "comes from lot N.° 15, under
        'shan', (to correct); so it's the wrong rhyme."}
}
{
}

\Verse{70}
{
    \zA{}
}
{
    \eA{This rambling talk made them involuntarily burst out laughing. "In very
        fact," Pao-sh'ai laughed, "she's under a poetical frenzy, and it's all P'in
        Erh who has incited her." "The holy man says," Tai-yü smilingly rejoined,
        "that 'one must not be weary of exhorting people'; and if she comes, time
        and again, to ask me this and that how can I possibly not tell her?" "Let's
        take her to Miss Quarta's rooms," Li Wan smiled, "and if we could coax her
        to look at the painting, and bring her to her senses, it will be well."
        Speaking the while, she actually walked out of the room, and laying hold of
        her, she brought her through the Lotus Fragrance arbour to the bank of Warm
        Fragrance. Hsi Ch'un was tired and languid, and was lying on the window,
        having a midday siesta.}
}
{
}

\Verse{71}
{
    \zA{}
}
{
    \eA{The painting was resting against the partition-wall, and was screened with a
        gauze cover. With one voice, they roused Hsi Ch'un, and raising the gauze
        cover to contemplate her work, they saw that three tenths of it had already
        been accomplished. But their attention was attracted by the representation
        of several beautiful girls, inserted in the picture, so pointing at Hsiang
        Ling: "Every one who can write verses is to be put here," they said, "so be
        quick and learn." But while conversing, they played and laughed for a time,
        after which, each went her own way.}
}
{
}

\Verse{72}
{
    \zA{}
}
{
    \eA{Hsiang Ling was meanwhile preoccupied about her verses, so, when evening
        came, she sat facing the lamp absorbed in thought. And the third watch
        struck before she got to bed. But her eyes were so wide awake, that it was
        only after the fifth watch had come and gone, that she, at length, felt
        drowsy and fell fast asleep. Presently, the day dawned, and Pao-ch'ai woke
        up; but, when she lent an ear, she discovered (Hsiang Ling) in a sound
        sleep. "She has racked her brains the whole night long," she pondered. "I
        wonder, however, whether she has succeeded in finishing her task. She must
        be tired now, so I won't disturb her." But in the midst of her cogitations,
        she heard Hsiang Ling laugh and exclaim in her sleep: "I've got it. It
        cannot be that this stanza too won't be worth anything."}
}
{
}

\Verse{73}
{
    \zA{}
}
{
    \eA{"How sad and ridiculous!" Pao-ch'ai soliloquised with a smile. And, calling
        her by name, she woke her up. "What have you got?" she asked. "With that
        firmness of purpose of yours, you could even become a spirit! But before you
        can learn how to write poetry, you'll be getting some illness." Chiding her
        the while, she combed her hair and washed; and, this done, she repaired,
        along with her cousins, into dowager lady Chia's quarters. Hsiang Ling made,
        in fact, such desperate efforts to learn all about poetry that her system
        got quite out of order. But although she did not in the course of the day
        hit upon anything, she quite casually succeeded in her dreams in devising
        eight lines;}
}
{
}

\Verse{74}
{
    \zA{}
}
{
    \eA{so concluding her toilette and her ablutions, she hastily jotted them down,
        and betook herself into the Hsin Fang pavilion. Here she saw Li Wan and the
        whole bevy of young ladies, returning from Madame Wang's suite of
        apartments. Pao-ch'ai was in the act of telling them of the verses composed
        by Hsiang Ling, while asleep, and of the nonsense she had been talking, and
        every one of them was convulsed with laughter. But upon raising their heads,
        and perceiving that she was approaching, they vied with each other in
        pressing her to let them see her composition. But, reader, do you wish to
        know any further particulars? If you do; read those given in the next
        chapter.}
}
{
}
